% Compile from this directory with: pdflatex project_13.tex (twice).
\documentclass[11pt,a4paper]{article}
\usepackage[T1]{fontenc}
\usepackage[utf8]{inputenc}
\usepackage[margin=22mm]{geometry}
\usepackage{lmodern}
\usepackage{amsmath,amssymb}
\usepackage{enumitem}
\usepackage{xcolor}
\usepackage{microtype}
\usepackage{hyperref}
\definecolor{epflred}{HTML}{E3000B}
\hypersetup{colorlinks=true,urlcolor=epflred,linkcolor=epflred,
  pdftitle={ENG-654 Project 13: Repeated Contour Execution on custom\_3R}}
\setlength{\parindent}{0pt}
\setlength{\parskip}{6pt}
\setlist[itemize]{leftmargin=*,itemsep=5pt,topsep=4pt}
\urlstyle{tt}

\begin{document}
{\small\textbf{EPFL \textbar\ ENG-654}\hfill Kinematics-Grounded Motion Planning for Robots}
\vspace{5mm}

{\color{epflred}\Large\bfseries Project 13}

{\LARGE\bfseries Repeated Contour Execution on custom\_3R\par}

\textbf{Format:} Group of two students; simulation-based project.\\
\textbf{Course foundations:} Lectures 3, 5--7; Exercise 3.

\section*{Motivation}
A closed workspace path need not return a cuspidal robot to its starting
posture. The next traversal then begins from a different IK and may behave
differently. Repeatability is therefore a property of the lifted joint motion,
not simply the visible closure of a Cartesian contour.

\section*{Task}
\textbf{Overall objectives.} Track repeated position contours on the exact custom\_3R robot and explain
joint-space closure, posture changes, and failed continuation using reduced
singularity maps and the workspace IK distribution.

\begin{itemize}
  \item \textbf{Prepare the model and maps.} Validate the exact orthogonal,
non-intersecting-axis custom\_3R kinematics. Plot $\det J_p$ and its zeros in
$(q_2,q_3)$, map the corresponding critical curves into $(\rho,z)$, and identify
the two-/four-IK workspace regions with representative IK overlays.

  \item \textbf{Choose two closed position curves.} Use an instructor-supplied
curve exhibiting nontrivial continuation and one nearby comparison curve.
Alternatively construct candidates with the Lecture 6 exact-model lab.
Verify closure in three-dimensional Cartesian coordinates; the Lecture 7
idealized square must be revalidated before transferring it to this robot.

  \item \textbf{Lift successive laps.} Enumerate the starting IKs and continue
each admissible one for up to three laps. Start every new lap at the actual
configuration reached by the previous lap, without resetting the seed.
Track solutions by continuity and stop or refine when a regular continuation
can no longer be recovered.

  \item \textbf{Classify the observed outcomes.} Distinguish joint-space closure,
return to a different IK of the same target, and termination. At each lap
boundary, match the reached configuration to the starting-pose IK set modulo
joint periodicity, and display the observed mapping. Do not assume every
curve defines a permutation if some lifts terminate.

  \item \textbf{Explain events globally.} Overlay all lifted paths on the reduced
singularity map and relate events to crossings of workspace critical curves
and changes in IK count. A critical-value crossing does not force every
possible lift to be singular; evaluate the Jacobian of the selected lift.
Refine the decisive intervals.
\end{itemize}

\newpage
\section*{Specifics and scope}
Use two loops, up to four initial mathematical arm IKs, and at most
three laps per initial solution. Separate joint-limit feasibility from
unconstrained topology. Reuse the analytical IK and continuation tools;
discovering a new cusp or proving a general repeatability theorem is outside
scope. A predicted outcome that fails numerical validation must be reported
rather than imposed.

Use the position Jacobian $J_p$ for the 3R task.
Confirm that the meridian reduction is valid for the selected model and
operational point. Count distinct real IKs modulo periodicity, then filter
joint limits separately. Identify two-/four-IK regions where present and
explain any absent count rather than forcing both types into every map. Show both determinant values and the zero contour,
and distinguish sampled connectivity from a proof. Refine the decisive
region instead of claiming completeness from a coarse grid.

\section*{Expected outcome}
Understand why repeated Cartesian motions can have different joint-space
outcomes. Deliver lap-by-lap configuration records, IK endpoint associations,
determinant and workspace-count plots, and an explanation of every observed
closure, posture change, or termination.

Submit reproducible code, model parameters and test data, the required plots,
a concise 5--6-page report, and a short simulation demonstration. Explain
both successful cases and failures; conclusions must follow the numerical
and geometric evidence rather than an expected outcome.

\section*{Resources}
The exact custom\_3R is an orthogonal 3R robot with non-intersecting
consecutive joint axes. Its URDF is
\path{lectures_main/assets/models/custom_3R/custom_3R_new.urdf}; STL meshes are
in the same directory. Use the lecture website to inspect axes and frames,
verify D--H parameters, plot IKs, and animate configurations.
Lecture 7 provides the repeated-loop concept, while Lecture 6 contains an
exact custom\_3R motion example. Two suitable starting curves or seed data
will be supplied if the existing demonstration cannot be exported directly.

Use the lecture website to verify D--H parameters, visualize IKs and
singularities, and simulate paths. All listed paths are relative to the
repository root. Start the website following \path{lectures_main/README.md}.

\section*{Workload and collaboration}
Budget approximately 60 hours per student (120 person-hours per pair)
for this 2-ECTS project, following EPFL's 30-hours-per-credit convention.\footnote{\href{https://www.epfl.ch/education/teaching/teaching-guide-2/getting-started/design-a-course_1/}{EPFL: Design a Course, workload guidance}.}
Deduct any lectures or exercises included in those credits. Allocate roughly
20 team-hours to model validation, 50 to the core work, 30 to experiments,
and 20 to reporting. Reuse the specified IK and simulations. Share validation
and interpretation even if modelling and implementation are divided between
students. Hardware execution is outside scope.

\end{document}
